\documentclass[lettersize,journal]{IEEEtran}
\usepackage{array}
\usepackage[caption=true,font=footnotesize,labelfont=rm,textfont=rm]{subfig} \usepackage{textcomp}
\usepackage{stfloats}
\usepackage{url}
\usepackage{verbatim}
\usepackage{graphicx}
\hyphenation{op-tical net-works semi-conduc-tor IEEE-Xplore}
\def\BibTeX{{\rm B\kern-.05em{\sc i\kern-.025em b}\kern-.08em
    T\kern-.1667em\lower.7ex\hbox{E}\kern-.125emX}}
\usepackage{balance}

\usepackage{amssymb,amsfonts} \usepackage[tbtags]{amsmath}

\usepackage[ruled,linesnumbered]{algorithm2e}
\usepackage{color} \usepackage{booktabs} \usepackage[thmmarks,amsmath]{ntheorem} \usepackage{setspace} \usepackage{hyperref}
\usepackage{cite}
\usepackage{soul}
\soulregister{\cite}7 \captionsetup{font={footnotesize}}
{
    \theoremstyle{nonumberplain}     \theoremheaderfont{\bfseries\itshape}     \theorembodyfont{\normalfont}     \theoremseparator{:}     \theoremsymbol{\ensuremath{\blacksquare}}     \newtheorem{proof}{Proof}
}
{
    \theoremstyle{plain}     \theoremheaderfont{\bfseries\itshape}     \theorembodyfont{\normalfont}     \theoremseparator{:}     \theoremsymbol{\ensuremath{\blacksquare}}     \newtheorem{remark}{Remark}
    \newtheorem{definition}{Definition}
    \newtheorem{lemma}{Lemma}
    \newtheorem{theorem}{Theorem}
}

\newcommand{\algref}[1]{\textit{\textbf{Algorithm \ref{#1}}}}
\newcommand{\theoref}[1]{\textit{\textbf{Theorem \ref{#1}}}}
\newcommand{\lemmref}[1]{\textit{\textbf{Lemma \ref{#1}}}}
\newcommand{\defref}[1]{\textit{\textbf{Definition \ref{#1}}}}
\newcommand{\figref}[1]{Fig. \ref{#1}}
\pdfoutput=1
\begin{document}
\title{Guided-DiffFNO: Policy Optimization via Critic-Guided Diffusion Fourier Neural Operators in HVAC Control
}

\author{Wei Zou,Yu Sun, Haibo Zhou,~\IEEEmembership{Fellow,~IEEE}}

\maketitle

\begin{abstract}
Building HVAC control presents a complex optimization challenge, requiring the simultaneous balancing of energy efficiency, thermal comfort, and multi-zone coordination. Conventional reinforcement learning (RL) approaches often struggle in this highly coupled environment, frequently producing unstable, high-frequency control actions due to their reliance on unimodal distributions and point-wise function approximators. To address these limitations, we propose Guided-DiffFNO, a generative control framework that combines conditional diffusion policies with spectral inductive biases. Instead of treating thermal zones independently, our method utilizes a Fourier Neural Operator (FNO) backbone to capture cross-zone spatial dependencies in the action space, which helps suppress undesirable high-frequency control jitter. Furthermore, we integrate an inference-time critic guidance mechanism, injecting value gradients into the reverse diffusion process to steer multi-zone actions toward safer, higher-reward regions. Evaluated in the physics-based BEAR simulation environment, Guided-DiffFNO attains a favorable energy--comfort trade-off compared to the RL baselines considered here. The integration of generative action modeling, spectral structure, and value-aware refinement also yields smoother and more physically consistent control trajectories, suggesting a promising direction for safe and efficient building control.
\end{abstract}

\begin{IEEEkeywords}
Building HVAC control, diffusion policy, Fourier neural operator, critic guidance, reinforcement learning, thermal comfort, energy-efficient control.
\end{IEEEkeywords}

\section{Introduction}

Buildings account for a substantial share of global energy consumption, and heating, ventilation, and air-conditioning (HVAC) systems are among the dominant contributors to building operational energy use \cite{perez2008review,cao2016building,cuce2026role}. Improving HVAC control is therefore of both practical and societal importance, as it directly affects energy efficiency, occupant comfort, and equipment operating stability. In modern commercial buildings, this problem is particularly challenging because thermal dynamics are strongly influenced by weather conditions, internal loads, occupancy patterns, and cross-zone heat transfer. As a result, high-performance HVAC control requires not only accurate local regulation, but also coordinated decision-making across multiple coupled thermal zones.

Classical control strategies, such as rule-based control, PID regulation, and model predictive control (MPC), have been widely studied for HVAC management. Rule-based and PID methods are easy to deploy, but their performance is often limited under highly dynamic operating conditions. MPC can explicitly account for future disturbances and physical constraints, and is frequently regarded as a strong baseline for building control. However, its practical effectiveness depends on model fidelity, online optimization efficiency, and computational tractability, all of which become increasingly difficult to maintain in high-dimensional multi-zone settings \cite{drgovna2020all}. These limitations have motivated growing interest in reinforcement learning (RL), which offers a data-driven alternative for long-horizon sequential decision-making without requiring exact analytical models of the environment.

Recent RL-based HVAC studies, comparative evaluations, and reviews have demonstrated the promise of deep actor--critic methods for improving energy performance while preserving thermal comfort \cite{wei2017deep,biemann2021experimental,mason2019review,alsayed2024hvacreview}. Nevertheless, standard RL policies still face several limitations in multi-zone building control. First, many continuous-control methods rely on deterministic mappings or unimodal Gaussian action distributions, which may be insufficient for representing the complex and potentially multi-modal structure of optimal HVAC decisions. Second, commonly used policy backbones, such as multilayer perceptrons (MLPs), treat action coordinates largely independently and thus fail to explicitly encode the coupled structure of zone-wise control signals. Third, these methods often produce unstable or highly oscillatory actions, especially when learning from sparse or delayed performance feedback. In building systems, such oscillatory behavior is undesirable because it may lead to unnecessary energy expenditure, comfort degradation, and poor physical consistency.

Diffusion models have recently emerged as a promising class of generative models for decision-making and control \cite{chi2025diffusion,janner2022planning}. By formulating action generation as an iterative denoising process, diffusion policies provide a richer modeling framework than conventional parametric policies and can represent complex action distributions without imposing restrictive unimodal assumptions. This property is particularly attractive for HVAC control, where multiple feasible action patterns may exist under the same thermal state due to coupled dynamics and disturbance uncertainty. However, directly applying a generic diffusion policy to building control is still insufficient. On the one hand, standard denoisers do not explicitly account for the structured coupling among zone-level control variables. On the other hand, unguided diffusion sampling remains primarily distribution-driven and does not necessarily bias generation toward actions with high long-term value or improved comfort compliance.

To address these issues, this paper proposes \emph{Guided-DiffFNO}, a generative control framework tailored for multi-zone building energy management. While preserving the expressive advantage of conditional diffusion models, we introduce two critical structural enhancements. First, we replace the standard MLP denoiser with an FNO-based architecture that operates explicitly on the action-dimension axis. This effectively imposes a spectral inductive bias, promoting globally coordinated, low-frequency control patterns across thermal zones. Second, we incorporate an inference-time critic guidance mechanism, utilizing the gradient of a learned double-critic value function to progressively steer the reverse denoising trajectory toward higher-value and lower-violation action regions. 

To stabilize the learning of this generative policy, we further develop a unified training framework that combines diffusion-consistent behavior cloning (utilizing expert MPC demonstrations) with conservative value-driven policy improvement. 

The main technical contributions of this work are summarized as follows:
\begin{itemize}
    \item \textbf{Spectral Action Modeling via DiffFNO:} We formulate a denoising architecture by deploying a Fourier Neural Operator across the \emph{action-dimension axis}. By transforming multi-zone action coordinates into the frequency domain, this backbone introduces a low-frequency spatial inductive bias that helps suppress physically inconsistent high-frequency control chatter across adjacent thermal zones.
    \item \textbf{Inference-Time Critic Guidance:} We design an inference-time guidance mechanism for continuous building control. By injecting gradients from a conservative double-Q value function into the reverse diffusion sampling steps, we steer the generated action manifold toward higher-return and safety-compliant regions while preserving the original diffusion prior.
    \item \textbf{Expert-Regularized Generative RL Paradigm:} We develop a hybrid actor--critic training scheme tailored for diffusion policies. This framework seamlessly unifies diffusion-consistent behavior cloning (from MPC experts) to stabilize early learning with value-driven policy gradient optimization for asymptotic performance improvement.
    \item \textbf{Empirical Validation of Physical Consistency:} Evaluations on the physics-based BEAR benchmark show that Guided-DiffFNO achieves a strong energy--comfort trade-off while also generating smoother, more globally coordinated control trajectories than the unimodal RL and unguided diffusion baselines considered in this study.
\end{itemize}

The remainder of this paper is organized as follows. Section\ref{sec:Related Work} reviews related work. Section\ref{sec:System Model} introduces the system model. Section\ref{sec:Problem Formulation and Solution} presents the proposed Guided-DiffFNO method. Section\ref{sec:Experimental Setup} describes the experimental setup. Section\ref{sec:results} discusses the experimental results, and Section\ref{sec:conclusion} concludes the paper.


\section{Related Work}\label{sec:Related Work}

\subsection{HVAC Control and Joint Multi-Zone Decision Making}
Building HVAC control has long been dominated by MPC because it can explicitly optimize energy and comfort under constraints \cite{drgovna2020all}. Its practical performance, however, still depends on model fidelity and repeated online optimization, both of which become more difficult in large multi-zone buildings. RL offers a model-free alternative and has shown promising results for HVAC operation in both single-zone and multi-zone settings \cite{wei2017deep,biemann2021experimental,mason2019review,alsayed2024hvacreview,ding2024multi}. Yet many existing methods still rely on deterministic or unimodal policies, per-zone decompositions, or other local action parameterizations, even when coordination among adjacent zones is important \cite{vazquez2020marlisa,yu2021multiagent,xue2025multi,cui2024energy,fang2023cross,dai2023reinforcement}. This leaves a gap for policies that directly generate the full joint action vector while respecting cross-zone coupling.

This gap is especially relevant in buildings where the optimization target is inherently system-level rather than zone-local. A controller that reduces the temperature error of one room by inducing aggressive neighboring-zone actuation may appear locally beneficial under a decomposed formulation, yet still be undesirable once energy use, comfort risk, and cross-zone interactions are considered jointly. Multi-agent or decentralized formulations can reduce action dimensionality, but they often shift the burden of coordination to reward shaping, communication design, or post hoc arbitration among local agents \cite{vazquez2020marlisa,yu2021multiagent,xue2025multi}. In contrast, our focus is on centralized \emph{joint-action generation}: the policy should model the full multi-zone control vector as a structured object and learn coordination directly in the action representation itself.

\subsection{Diffusion Policies for Continuous Control}
Diffusion models provide a more expressive alternative to conventional continuous-control policies because they can model complex action distributions through iterative denoising. Diffusion Policy \cite{chi2025diffusion} demonstrated strong performance on high-dimensional control, while Diffuser \cite{janner2022planning} and Diffusion-QL \cite{wang2023diffusion} extended diffusion ideas to RL and planning. To the best of our knowledge, such policies have not been systematically studied for building HVAC control. Moreover, existing diffusion policies typically use MLP denoisers, which do not explicitly encode the spatial coupling of multi-zone actions.

Most existing diffusion-control studies also target domains whose structure differs substantially from building control. In robotics, diffusion models are often used to generate temporally extended action chunks or visuomotor trajectories, and the main modeling challenge is multimodality over time \cite{chi2025diffusion,janner2022planning}. In HVAC control, however, an equally important challenge lies in the \emph{spatial} structure of the action itself: each dimension corresponds to a thermal zone embedded in a coupled physical system. A diffusion formulation is attractive here not only because it relaxes unimodal policy assumptions, but also because its iterative denoising process provides a natural interface for injecting structure and value information into the generated joint action.

\subsection{Structural Priors and Guided Diffusion}
FNOs \cite{li2021fourier} provide a natural structural prior when long-range coupling matters, and their tendency to emphasize dominant low-frequency components can be beneficial when coordinated control should avoid cross-zone chatter \cite{qin2024toward,liu2024mitigating}. We therefore use FNO not to claim that HVAC actions obey a PDE solution operator, but to bias the denoiser toward globally coherent joint actions under thermal coupling \cite{goyal2012method,wen2022u}. In parallel, guided diffusion has shown that inference-time gradients can steer generation without retraining the denoiser \cite{dhariwal2021diffusion,ho2021classifier}. RL-oriented variants have begun to use policy or value signals for this purpose \cite{jackson2024policy,fang2025diffusion}, but they often rely on additional guidance networks or training-time modifications. Our method instead combines an FNO denoiser with inference-time guidance from the critic already required by the actor--critic framework.

This positioning differentiates our work from two adjacent lines of research. First, neural operators in building applications are commonly used as fast surrogates for thermal dynamics or PDE-governed fields; in our setting, by contrast, the operator acts inside the \emph{policy} and regularizes the geometry of the generated action vector rather than approximating the environment transition. Second, many guided-diffusion methods in RL improve generation quality by introducing auxiliary classifiers, reward models, or preference models. Our design choice is intentionally lighter: we reuse the conservative critic already needed for off-policy learning, and we invoke it only during reverse sampling so that the base diffusion policy remains unchanged. This separation also makes the method easier to interpret experimentally, because the contributions of structural prior and value-guided refinement can be studied independently.


\section{System Model}\label{sec:System Model}

\begin{figure}[!t]
    \captionsetup{width=\linewidth}
    \centerline{\includegraphics[width=\linewidth]{figures/system model.pdf}}
    \caption{System Schematic of the Multi-Zone Building HVAC Control Problem}
    \label{fig:system_model}
\end{figure}

We formulate multi-zone building HVAC control as a discounted Markov decision process (MDP)
\begin{equation}
    \mathcal{M} = (\mathcal{S}, \mathcal{A}, \mathcal{P}, r, \gamma),
\end{equation}
where $\mathcal{S}$ and $\mathcal{A}$ denote the state and action spaces, $\mathcal{P}$ is the transition kernel induced by the building simulator, $r$ is the reward function, and $\gamma \in (0,1]$ is the discount factor. At each control step $t$, the controller observes the current building state $\mathbf{s}_t \in \mathcal{S}$ and outputs a continuous joint control action $\mathbf{a}_t \in \mathcal{A}$ for all thermal zones. The control objective is to reduce HVAC energy consumption while maintaining indoor thermal comfort under coupled building dynamics and external disturbances.

\subsection{Building Thermal Environment}
We consider an $N$-zone building modeled in the physics-based BEAR simulator\cite{zhang2023bear}. The thermal evolution of the building is jointly affected by HVAC actuation, outdoor weather conditions, solar irradiance, ground temperature, and internal heat gains. Let $\mathbf{s}_t$ denote the system state, $\mathbf{a}_t$ the control action, and $\boldsymbol{\xi}_t$ the exogenous disturbance realization at time step $t$. The building dynamics can be written in compact form as
\begin{equation}
    \mathbf{s}_{t+1} = f(\mathbf{s}_t, \mathbf{a}_t, \boldsymbol{\xi}_t),
    \label{eq:dynamics}
\end{equation}
where $f(\cdot)$ is the transition function implemented by the underlying simulator. In this work, the simulator serves as the environment for policy learning, and the controller does not assume access to an explicit analytical form of $f(\cdot)$.

A key characteristic of the problem is that the dynamics are \emph{cross-zone coupled}: the temperature evolution of each zone depends not only on its local actuation but also on neighboring-zone interactions and shared disturbances. Therefore, the control problem should be treated as \emph{joint action generation} over all zones rather than a collection of independent single-zone decisions.

\subsection{State Space}
The state vector contains both thermal states and disturbance-related variables. For an $N$-zone building, the state at time step $t$ is defined as
\begin{equation}
    \mathbf{s}_t =
    \left[
    \mathbf{T}_t^{\mathrm{zone}},
    T_t^{\mathrm{out}},
    \mathbf{G}_t,
    T_t^{\mathrm{ground}},
    \mathbf{L}_t
    \right]
    \in \mathbb{R}^{3N+2},
    \label{eq:state}
\end{equation}
where $\mathbf{T}_t^{\mathrm{zone}} \in \mathbb{R}^N$ denotes the indoor temperatures of all zones, $T_t^{\mathrm{out}} \in \mathbb{R}$ is the outdoor air temperature, $\mathbf{G}_t \in \mathbb{R}^N$ represents zone-associated solar irradiance features, $T_t^{\mathrm{ground}} \in \mathbb{R}$ is the ground temperature, and $\mathbf{L}_t \in \mathbb{R}^N$ denotes the internal thermal loads of all zones, such as occupancy- and equipment-related heat gains. Hence, the total state dimension is $3N+2$.

This state representation preserves the coupled thermal structure of the multi-zone building and provides sufficient information for control-oriented decision making.

\subsection{Action Space}
The control action is defined as a continuous HVAC command for each zone:
\begin{equation}
    \mathbf{a}_t =
    [a_{t,1}, a_{t,2}, \dots, a_{t,N}]
    \in [-1,1]^N.
    \label{eq:action}
\end{equation}
Each action component corresponds to the normalized HVAC actuation intensity of one zone. Negative values indicate cooling and positive values indicate heating. The action is bounded to ensure physically meaningful actuation and stable interaction with the simulator.

Importantly, $\mathbf{a}_t$ is a \emph{joint} control vector over all zones. This formulation is consistent with the coupled dynamics in \eqref{eq:dynamics} and emphasizes that coordinated multi-zone control is required for good performance.

Given the normalized action, the instantaneous HVAC power is computed from the zone-wise actuation map and the maximum equipment power. Let $m_i \ge 0$ denote the actuation scaling coefficient of zone $i$, obtained from the simulator, and let $P_{\max}$ denote the maximum unit power in watts. Then the total HVAC power at time step $t$ is
\begin{equation}
    P_t^{\mathrm{hvac}}
    =
    \frac{P_{\max}}{1000}
    \sum_{i=1}^{N} m_i |a_{t,i}|,
    \qquad [\mathrm{kW}],
    \label{eq:power}
\end{equation}
and the corresponding energy consumption over one control interval $\Delta t$ seconds is
\begin{equation}
    E_t^{\mathrm{hvac}}
    =
    P_t^{\mathrm{hvac}} \frac{\Delta t}{3600},
    \qquad [\mathrm{kWh}].
    \label{eq:energy}
\end{equation}

\subsection{Thermal Comfort and Reward Design}
The control objective must balance energy efficiency and occupant comfort. Let $T^{\mathrm{ref}}$ denote the target indoor temperature and $\delta$ the allowed comfort tolerance following ASHRAE Standard 55 \cite{standard2023thermal}. For each zone $i$, the temperature deviation is
\begin{equation}
    d_{t,i} = \left| T_{t,i}^{\mathrm{zone}} - T^{\mathrm{ref}} \right|.
\end{equation}
Accordingly, the mean comfort deviation and the number of comfort violations at time step $t$ are defined as
\begin{equation}
    D_t^{\mathrm{mean}} = \frac{1}{N} \sum_{i=1}^{N} d_{t,i},
    \label{eq:comfort_mean}
\end{equation}
\begin{equation}
    V_t = \sum_{i=1}^{N} \mathbb{I}(d_{t,i} > \delta),
    \label{eq:violation}
\end{equation}
where $\mathbb{I}(\cdot)$ is the indicator function. In addition, the maximum temperature deviation is
\begin{equation}
    D_t^{\mathrm{max}} = \max_{1 \le i \le N} d_{t,i},
    \label{eq:comfort_max}
\end{equation}
which is used as an auxiliary evaluation metric.

The simulator provides a base reward $r_t^{\mathrm{env}}$ that reflects the trade-off between energy usage and temperature regulation. To further discourage unsafe thermal deviations, we add an explicit comfort-violation penalty and apply reward scaling for stable learning. The final reward is defined as
\begin{equation}
    r_t = \rho \left( r_t^{\mathrm{env}} - \lambda V_t \right),
    \label{eq:reward}
\end{equation}
where $\lambda > 0$ is the violation penalty coefficient and $\rho > 0$ is a reward scaling factor. This reward design encourages the controller to reduce HVAC energy use while keeping all zone temperatures within the desired comfort range.

\subsection{Control Objective}
The goal is to learn a stochastic control policy $\pi(\mathbf{a}_t \mid \mathbf{s}_t)$ that maximizes the expected cumulative discounted return over an episode of horizon $T$:
\begin{equation}
    J(\pi)
    =
    \mathbb{E}_{\pi}
    \left[
    \sum_{t=0}^{T-1} \gamma^t r_t
    \right].
    \label{eq:objective}
\end{equation}
Equivalently, the learned controller should optimize the long-term trade-off among HVAC energy consumption, thermal comfort, and comfort-violation risk under disturbance uncertainty and cross-zone thermal coupling.

The formulation above also clarifies the main challenge of the problem: because the transition dynamics are coupled across zones, a high-quality policy must generate coordinated joint actions that are globally consistent with the building thermal field, rather than treating each zone as an isolated control channel.


\section{Problem Formulation and Solution}
\label{sec:Problem Formulation and Solution}
\label{sec:method}

Building on the MDP formulation in Section~\ref{sec:System Model}, 
we highlight three distinctive challenges that make multi-zone HVAC 
control particularly demanding for standard continuous-control methods.

First, the state-action relationship is highly nonlinear because zone temperatures are jointly affected by cross-zone thermal coupling, weather disturbances, solar irradiance, and internal heat gains. Second, the action vector has an inherent spatial structure: each dimension corresponds to one thermal zone, and the action should therefore be generated as a coordinated joint control vector rather than as a set of independent commands. Third, actions that appear locally feasible may still be suboptimal from the perspective of long-term value, especially when the building operates near thermal comfort boundaries.

These characteristics motivate a policy class that is distributionally expressive, structurally aware, and amenable to value-based refinement during inference. To address these challenges, we propose Guided-DiffFNO, a diffusion-policy-based framework that combines conditional denoising, spectral action modeling, and critic guidance during reverse sampling.

\subsection{Conditional Diffusion Policy for Continuous HVAC Control}

Instead of directly regressing a deterministic action from the current state, we model policy generation as a conditional denoising diffusion process over the action space \cite{ho2020denoising}. Let $\mathbf{x}_0 \in \mathbb{R}^N$ denote the clean action and $\mathbf{x}_k$ its noisy version at diffusion step $k \in \{1,\dots,K\}$. Under the standard VP diffusion parameterization,
\begin{equation}
\mathbf{x}_k
=
\sqrt{\bar{\alpha}_k}\mathbf{x}_0
+
\sqrt{1-\bar{\alpha}_k}\boldsymbol{\epsilon},
\qquad
\boldsymbol{\epsilon}\sim\mathcal{N}(\mathbf{0},\mathbf{I}).
\label{eq:noisy_action}
\end{equation}
where $\alpha_k = 1-\beta_k$ and $\bar{\alpha}_k = \prod_{i=1}^{k}\alpha_i$. In our implementation, the denoiser is parameterized to \emph{directly predict the clean action}:
\begin{equation}
\hat{\mathbf{x}}_0 = f_{\theta}(\mathbf{x}_k,k,\mathbf{s}_t),
\label{eq:x0_prediction}
\end{equation}
where $f_\theta(\cdot)$ takes the noisy action, the diffusion step, and the current state as input. The reverse posterior then follows the standard DDPM form of Ho \emph{et al.} \cite{ho2020denoising}; under the $\mathbf{x}_0$-prediction parameterization, its mean is
\begin{equation}
\boldsymbol{\mu}_k(\hat{\mathbf{x}}_0,\mathbf{x}_k)
=
c_1(k)\hat{\mathbf{x}}_0 + c_2(k)\mathbf{x}_k,
\label{eq:posterior_mean}
\end{equation}
where $c_1(k)$ and $c_2(k)$ are the standard DDPM coefficients determined by the VP schedule. The posterior variance is determined by the same schedule. Starting from $\mathbf{x}_K \sim \mathcal{N}(\mathbf{0},\mathbf{I})$, the policy iteratively samples backward to obtain the final action, which is clipped into the feasible range $[-1,1]^N$. Compared with unimodal Gaussian policies, this formulation offers a more expressive action model and a natural interface for structural priors and value-based guidance.

Fig.~\ref{fig:guided_difffno_framework} summarizes the overall architecture and training/inference workflow of the proposed Guided-DiffFNO framework.

\begin{figure*}[!t]
    \centering
    \includegraphics[width=\textwidth]{figures/guided_difffno_framework.pdf}
    \caption{Overall architecture of Guided-DiffFNO. Top: critic-guided reverse diffusion inference with the DiffFNO denoiser. Middle: DiffFNO backbone operating on the action-dimension axis. Bottom: off-policy training with critic learning and hybrid actor optimization that combines behavior cloning and policy-gradient objectives.}
    \label{fig:guided_difffno_framework}
\end{figure*}


\subsection{DiffFNO Backbone}

A key limitation of conventional denoisers based on multilayer perceptrons is that they treat the action vector as a generic coordinate array and therefore do not explicitly exploit the structured coupling among zone-wise control signals. In multi-zone HVAC control, however, the action is not merely a collection of independent commands. Rather, it forms a coordinated control field shaped by coupled thermal interactions across zones. To encode this structure, we replace the standard denoising backbone with a Fourier neural operator (FNO), yielding the proposed \emph{DiffFNO} architecture.

Given a noisy action $\mathbf{x}_k \in \mathbb{R}^N$, the current diffusion step $k$, and the state $\mathbf{s}_t$, DiffFNO first constructs a state embedding and a diffusion-step embedding,
\begin{equation}
\mathbf{h}_s=\phi_s(\mathbf{s}_t),
\qquad
\mathbf{h}_k=\phi_k(k),
\end{equation}
which are concatenated and projected into a latent conditioning vector
\begin{equation}
\mathbf{h}_c = \phi_c([\mathbf{h}_s,\mathbf{h}_k]) \in \mathbb{R}^{w},
\end{equation}
where $w$ denotes the latent channel width. The noisy action is then lifted into a one-dimensional latent feature map,
\begin{equation}
\mathbf{Y}^{(0)} = \phi_{\mathrm{in}}(\mathbf{x}_k) \in \mathbb{R}^{w \times N},
\end{equation}
and the conditioning vector is injected into the latent representation.

For each spectral layer $\ell$, the latent feature is updated as
\begin{equation}
\mathbf{Y}^{(\ell+1)}
=
\sigma
\left(
\mathcal{K}_{\mathrm{spec}}^{(\ell)}(\mathbf{Y}^{(\ell)})
+
\mathcal{K}_{\mathrm{pw}}^{(\ell)}(\mathbf{Y}^{(\ell)})
+
\mathbf{h}_c
\right),
\label{eq:spectral_layer}
\end{equation}
where $\mathcal{K}_{\mathrm{spec}}^{(\ell)}$ denotes a one-dimensional spectral convolution, $\mathcal{K}_{\mathrm{pw}}^{(\ell)}$ denotes a pointwise $1\times 1$ convolution, and $\sigma(\cdot)$ is the activation function.

To make the spectral operation explicit, let $\widehat{\mathbf{Y}}^{(\ell)}=\mathcal{F}(\mathbf{Y}^{(\ell)})$ denote the real-to-complex Fourier transform of the latent feature along the action-dimension axis. For the first $M$ retained modes, the spectral branch applies learnable complex linear mappings:
\begin{equation}
\widehat{\mathbf{Z}}^{(\ell)}_{:,m}
=
\mathbf{R}^{(\ell)}_{m}\widehat{\mathbf{Y}}^{(\ell)}_{:,m},
\qquad
m=1,\dots,M,
\label{eq:spectral_modes}
\end{equation}
where $\mathbf{R}^{(\ell)}_{m} \in \mathbb{C}^{w\times w}$ is the complex weight matrix associated with mode $m$. Higher modes are discarded:
\begin{equation}
\widehat{\mathbf{Z}}^{(\ell)}_{:,m}=\mathbf{0},
\qquad
m>M.
\label{eq:mode_truncation}
\end{equation}
The spectral output is then obtained through inverse Fourier transform:
\begin{equation}
\mathcal{K}_{\mathrm{spec}}^{(\ell)}(\mathbf{Y}^{(\ell)})
=
\mathcal{F}^{-1}\left(\widehat{\mathbf{Z}}^{(\ell)}\right).
\label{eq:spec_conv}
\end{equation}

It is important to emphasize that the FNO in this work operates along the \emph{action-dimension axis}, rather than along the temporal axis. Thus, it does not explicitly model temporal action trajectories. Instead, it imposes a spectral inductive bias directly on the multi-zone action vector, encouraging globally coordinated low-frequency control patterns and suppressing spatially inconsistent high-frequency oscillations across zones. This design is well aligned with the physical nature of multi-zone HVAC control, where mutually conflicting zone commands are generally undesirable.

After the spectral stack, the latent feature is projected back to the action space to produce the spectral denoising output
\begin{equation}
\hat{\mathbf{y}}_k = \phi_{\mathrm{out}}(\mathbf{Y}^{(L)}).
\end{equation}
To preserve local flexibility and avoid excessive over-smoothing, we further introduce a residual linear pathway
\begin{equation}
\mathbf{r}_k = \mathbf{W}_r \mathbf{x}_k,
\end{equation}
and define the final denoiser output as
\begin{equation}
f_{\theta}(\mathbf{x}_k,k,\mathbf{s}_t)
=
\hat{\mathbf{y}}_k + \mathrm{sigmoid}(g)\mathbf{r}_k,
\label{eq:final_denoiser}
\end{equation}
where $g$ is a learnable scalar gate. The residual pathway allows the model to retain local action variations useful for zone-specific regulation, while the spectral branch provides a global coordination prior. In this way, DiffFNO improves both expressiveness and structural consistency.

\subsection{Critic Guidance During Reverse Sampling}

Although the diffusion policy can generate physically consistent actions, consistency alone does not guarantee that the generated action is optimal with respect to long-term control performance. In particular, near thermal comfort boundaries, unguided reverse sampling may still converge to feasible but suboptimal actions. To alleviate this issue, we introduce critic guidance during reverse sampling.

Let $Q_{\phi_1}(\mathbf{s},\mathbf{a})$ and $Q_{\phi_2}(\mathbf{s},\mathbf{a})$ denote the two critic networks, and define the conservative value estimate as
\begin{equation}
Q_{\phi}^{\min}(\mathbf{s},\mathbf{a})
=
\min\left\{
Q_{\phi_1}(\mathbf{s},\mathbf{a}),
Q_{\phi_2}(\mathbf{s},\mathbf{a})
\right\}.
\label{eq:qmin}
\end{equation}
At reverse diffusion step $k$, the actor first reconstructs the clean action estimate $\hat{\mathbf{x}}_0 = f_\theta(\mathbf{x}_k,k,\mathbf{s}_t)$. The gradient of the conservative critic value with respect to the reconstructed action is then computed as
\begin{equation}
\mathbf{g}_k
=
\nabla_{\mathbf{a}} Q_{\phi}^{\min}(\mathbf{s}_t,\mathbf{a})
\big|_{\mathbf{a}=\hat{\mathbf{x}}_0}.
\label{eq:qgrad}
\end{equation}
This gradient is injected into the reverse chain by correcting the reconstructed action,
\begin{equation}
\tilde{\mathbf{x}}_0
=
\hat{\mathbf{x}}_0 + \eta \mathbf{g}_k,
\label{eq:guided_x0}
\end{equation}
where $\eta > 0$ is the guidance scale. In the implementation, this is equivalently written as
\begin{equation}
\tilde{\mathbf{x}}_0
=
\hat{\mathbf{x}}_0 - \eta \mathbf{d}_k,
\qquad
\mathbf{d}_k = -\mathbf{g}_k,
\label{eq:guided_impl}
\end{equation}
which corresponds to subtracting the negative value gradient. The corrected clean action is clipped into the feasible range and substituted into the posterior mean,
\begin{equation}
\boldsymbol{\mu}_k(\tilde{\mathbf{x}}_0,\mathbf{x}_k)
=
c_1(k)\tilde{\mathbf{x}}_0 + c_2(k)\mathbf{x}_k.
\label{eq:guided_mean}
\end{equation}
The reverse sampling process then proceeds accordingly.

This design has two advantages. First, critic guidance is integrated \emph{within} the reverse diffusion dynamics rather than applied as a post hoc correction after action generation. Consequently, value information influences the denoising trajectory throughout the sampling process. Second, because thermal comfort violations are already reflected in the reward, the learned critic implicitly discourages actions likely to induce unsafe temperature deviations. Therefore, critic guidance improves value alignment and empirically promotes comfort-compliant action generation. However, it should be interpreted as value-guided refinement rather than as a formal hard-safety guarantee.
At inference, the controller starts from Gaussian noise and, for each reverse step, applies the denoiser, evaluates the critic gradient in \eqref{eq:qgrad}, forms the guided reconstruction in \eqref{eq:guided_x0}, and samples the next state from the posterior with mean \eqref{eq:guided_mean}. The final action is clipped to $[-1,1]^N$ before execution.

\subsection{Critic Learning}

The value function is learned using a double-critic architecture. For each state-action pair $(\mathbf{s}_t,\mathbf{a}_t)$, the critic outputs two Q-values,
\begin{equation}
Q_{\phi_1}(\mathbf{s}_t,\mathbf{a}_t),
\qquad
Q_{\phi_2}(\mathbf{s}_t,\mathbf{a}_t).
\end{equation}
To reduce overestimation bias, target values are computed using the minimum of the two target critics\cite{fujimoto2018addressing}. Specifically, for the bootstrap state $\mathbf{s}_{t+n}$, the target actor generates the next action $\mathbf{a}'_{t+n}$, and the target Q-value is defined as
\begin{equation}
\begin{aligned}
Q_{\phi'}^{\min}(\mathbf{s}_{t+n},\mathbf{a}'_{t+n})
&= \min\Big\{ Q_{\phi_1'}(\mathbf{s}_{t+n},\mathbf{a}'_{t+n}), \\
&\quad\ Q_{\phi_2'}(\mathbf{s}_{t+n},\mathbf{a}'_{t+n}) \Big\}.
\end{aligned}
\end{equation}
Using $n$-step return estimation\cite{sutton2018reinforcement}, the target is written as
\begin{equation}
y_t^{(n)} = \sum_{i=0}^{n-1}\gamma^i r_{t+i} + \gamma^n Q_{\phi'}^{\min}(\mathbf{s}_{t+n},\mathbf{a}'_{t+n}).
\label{eq:nstep_target}
\end{equation}
The critic objective is then
\begin{equation}
\begin{aligned}
\mathcal{L}_{\mathrm{critic}} 
&= \mathbb{E}\left[\left(Q_{\phi_1}(\mathbf{s}_t,\mathbf{a}_t)-y_t^{(n)}\right)^2\right] \\
&\quad + \mathbb{E}\left[\left(Q_{\phi_2}(\mathbf{s}_t,\mathbf{a}_t)-y_t^{(n)}\right)^2\right].
\end{aligned}
\label{eq:critic_loss}
\end{equation}

The double-critic design is particularly important here because the critic is used not only for policy improvement during training but also as the guidance source during reverse diffusion inference. A more conservative value surface therefore benefits both training stability and inference quality.

\subsection{Actor Learning with Expert-Regularized Policy Optimization}

The actor is optimized through a combination of diffusion-based expert regularization and critic-driven policy improvement. Let $\pi_{\mathrm{exp}}$ denote the expert controller. In the final implementation, the expert label is aligned with the \emph{current} state, namely
\begin{equation}
\mathbf{a}_t^{\mathrm{exp}}
=
\pi_{\mathrm{exp}}(\mathbf{s}_t).
\label{eq:expert_alignment}
\end{equation}
This alignment avoids temporal mismatch between the supervised label and the policy input state.

Under the diffusion actor's clean-action prediction parameterization, a diffusion step $k$ is first sampled uniformly, and Gaussian noise is injected into the expert action:
\begin{equation}
\mathbf{x}_{k}^{\mathrm{exp}}
=
\sqrt{\bar{\alpha}_{k}}\,\mathbf{a}_t^{\mathrm{exp}}
+
\sqrt{1-\bar{\alpha}_{k}}\,\boldsymbol{\epsilon},
\qquad
\boldsymbol{\epsilon}\sim\mathcal{N}(\mathbf{0},\mathbf{I}).
\label{eq:expert_noisy}
\end{equation}
The denoiser is then trained to reconstruct the clean expert action:
\begin{equation}
\mathcal{L}_{\mathrm{BC}}
=
\mathbb{E}_{\mathbf{s}_t,\mathbf{a}_t^{\mathrm{exp}},k,\boldsymbol{\epsilon}}
\left[
\left\|
f_{\theta}(\mathbf{x}_{k}^{\mathrm{exp}},k,\mathbf{s}_t)
-
\mathbf{a}_t^{\mathrm{exp}}
\right\|_2^2
\right].
\label{eq:bc_loss}
\end{equation}
This objective is exactly consistent with the diffusion actor used in our implementation: the noisy expert action is the input, and the clean expert action is the training target.

In addition to expert regularization, the actor is optimized to maximize the conservative critic value. For a sampled state $\mathbf{s}_t$, the actor generates an action through reverse diffusion, denoted as $\pi_\theta(\mathbf{s}_t)$, and the policy-improvement term is defined as
\begin{equation}
\mathcal{L}_{\mathrm{PG}}
=
-
\mathbb{E}_{\mathbf{s}_t}
\left[
Q_{\phi}^{\min}(\mathbf{s}_t,\pi_\theta(\mathbf{s}_t))
\right].
\label{eq:pg_loss}
\end{equation}
The final actor objective is a weighted combination of the two terms, following the general spirit of behavior-regularized policy optimization in offline RL \cite{fujimoto2021minimalist}:
\begin{equation}
\mathcal{L}_{\mathrm{actor}}
=
\omega \mathcal{L}_{\mathrm{BC}}
+
(1-\omega)\mathcal{L}_{\mathrm{PG}},
\label{eq:actor_loss}
\end{equation}
where $\omega \in [0,1]$ is the behavior-cloning weight. In practice, $\omega$ is initialized with a relatively large value and then linearly decayed toward a smaller terminal value during training. Let $u$ denote the actor update step and $H$ denote the decay horizon. The schedule is
\begin{equation}
\omega(u)
=
\omega_0
+
\min\left(\frac{u}{H},1\right)(\omega_f-\omega_0),
\label{eq:bc_schedule}
\end{equation}
where $\omega_0$ and $\omega_f$ are the initial and final behavior-cloning weights, respectively. This schedule stabilizes early learning through expert supervision while gradually allowing the policy to move beyond imitation and optimize long-term control performance.

After each update, both the target actor and target critics are updated using Polyak averaging\cite{lillicrap2016continuous}:
\begin{equation}
\theta' \leftarrow \tau_{\mathrm{p}} \theta + (1-\tau_{\mathrm{p}})\theta',
\qquad
\phi' \leftarrow \tau_{\mathrm{p}} \phi + (1-\tau_{\mathrm{p}})\phi',
\label{eq:polyak}
\end{equation}
where $\tau_{\mathrm{p}}$ is the soft update rate.

\begin{algorithm}[t]
\caption{Training procedure of Guided-DiffFNO}
\label{alg:training_guided_difffno}
\KwIn{Replay buffer $\mathcal{B}$, actor $\pi_\theta$, target actor $\pi_{\theta'}$, critics $Q_{\phi_1}$ and $Q_{\phi_2}$, target critics $Q_{\phi_1'}$ and $Q_{\phi_2'}$, initial BC weight $\omega_0$, final BC weight $\omega_f$, decay horizon $H$.}
\KwOut{Trained actor $\pi_\theta$, critics $Q_{\phi_1}$ and $Q_{\phi_2}$.}

Initialize actor and critics, and copy them to the target networks\;
Initialize the behavior-cloning weight $\omega \leftarrow \omega_0$\;

\While{the stopping criterion is not met}{
    Collect transitions $(\mathbf{s}_t,\mathbf{a}_t,r_t,\mathbf{s}_{t+1},\mathrm{info}_t)$ and store them in $\mathcal{B}$\;
    Sample a minibatch from $\mathcal{B}$\;

    Generate target action $\mathbf{a}'_{t+n} = \pi_{\theta'}(\mathbf{s}_{t+n})$\;
    Compute conservative target value
    $Q_{\phi'}^{\min}(\mathbf{s}_{t+n},\mathbf{a}'_{t+n})
    =
    \min\{Q_{\phi_1'}(\mathbf{s}_{t+n},\mathbf{a}'_{t+n}),Q_{\phi_2'}(\mathbf{s}_{t+n},\mathbf{a}'_{t+n})\}$\;
    Compute the $n$-step target using \eqref{eq:nstep_target}\;
    Update the critics by minimizing $\mathcal{L}_{\mathrm{critic}}$\;

    Generate actor action $\tilde{\mathbf{a}}_t = \pi_\theta(\mathbf{s}_t)$\;
    Compute policy-improvement loss $\mathcal{L}_{\mathrm{PG}}$ using \eqref{eq:pg_loss}\;

    \eIf{expert action $\mathbf{a}_t^{\mathrm{exp}}$ exists in $\mathrm{info}_t$}{
        Sample diffusion step $k \sim \mathrm{Uniform}\{1,\dots,K\}$\;
        Sample Gaussian noise $\boldsymbol{\epsilon} \sim \mathcal{N}(\mathbf{0},\mathbf{I})$\;
        Construct noisy expert action using \eqref{eq:expert_noisy}\;
        Compute expert-regularization loss $\mathcal{L}_{\mathrm{BC}}$ using \eqref{eq:bc_loss}\;
        Set mixed actor loss
        $\mathcal{L}_{\mathrm{actor}} = \omega \mathcal{L}_{\mathrm{BC}} + (1-\omega)\mathcal{L}_{\mathrm{PG}}$\;
    }{
        Set $\mathcal{L}_{\mathrm{actor}} = \mathcal{L}_{\mathrm{PG}}$\;
    }

    Update actor parameters by minimizing $\mathcal{L}_{\mathrm{actor}}$\;
    Soft-update target networks using \eqref{eq:polyak}\;
    Update $\omega$ according to \eqref{eq:bc_schedule}\;
}
\Return{Trained actor and critics}\;
\end{algorithm}

\subsection{Inference Procedure}

During deployment, the current state $\mathbf{s}_t$ is passed to the diffusion actor, which initializes the reverse chain from Gaussian noise and iteratively applies DiffFNO denoising with optional critic guidance until the final action is obtained. The resulting controller combines generative action modeling, spectral structure learning over the joint action vector, and value-based refinement during sampling, enabling smoother and better coordinated HVAC control.




\section{Experimental Setup}
\label{sec:Experimental Setup}

\subsection{Benchmark and Environment}
We evaluate all methods on the BEAR \emph{OfficeSmall-Hot\_Dry} benchmark, corresponding to a six-zone office building under Tucson hot-dry weather. Unlike traditional EnergyPlus co-simulation setups, which suffer from prohibitive computational overhead and synchronization instability during millions of interactions, the BEAR environment provides a physics-grounded, vectorized simulator optimized for deep reinforcement learning. This setting allows us to evaluate generated control trajectories under rigorous thermal dynamics while preserving the computational efficiency required for diffusion-model training. The control interval is one hour, and each episode spans $168$ steps, i.e., one simulated week. Unless otherwise stated, all controllers share the same environment configuration: target temperature $26^\circ$C, comfort tolerance $\pm 1^\circ$C, energy weight $0.4$, temperature-regulation weight $0.6$, additional comfort-violation penalty coefficient $\lambda=10$, reward scaling factor $\rho=3.5\times 10^{-4}$, and maximum HVAC power $8000$~W. These settings ensure that all methods are evaluated under an identical energy--comfort objective.

\subsection{Compared Methods and Training Protocol}
We compare Guided-DiffFNO against four diffusion-family variants and three conventional baselines: DiffFNO w/o Guidance, DiffFNO w/o Residual \& Guidance, DiffFNO w/o Residual, Diffusion-MLP (i.e., Diffusion Policy with an MLP backbone), MPC, SAC\cite{haarnoja2018soft}, and SAC+MPC. For the diffusion-family methods, we use the same actor--critic training pipeline and expert-regularization schedule described in Section~\ref{sec:method}, changing only the component under study in each ablation. Expert demonstrations are generated by an MPC controller with planning horizon $3$.

To make the comparison as controlled as possible, all learning-based methods use the same state inputs, action bounds, reward function, replay setting, training horizon, and evaluation environment. In particular, the four diffusion-family variants differ only in the component being studied: whether the denoiser backbone is FNO or MLP, whether the residual pathway is present, and whether critic guidance is activated at inference. This design matters because diffusion methods are sensitive to auxiliary choices; without shared training and guidance settings, improvements could otherwise be attributed to retuning rather than to the proposed structural changes. For the conventional baselines, we preserve their native decision rules while evaluating them under the same simulator, weather profile, and comfort objective, so that the comparison isolates controller quality rather than differences in task definition.

The shared training and inference hyperparameters for the learning-based methods are summarized in Table~\ref{tab:training_setup}. We keep the guidance scale and the behavior-cloning decay schedule fixed across diffusion-family variants so that the ablations isolate architectural and guidance effects rather than retuning auxiliary settings.

\begin{table*}[t]
\centering
\caption{Shared hyperparameters for the learning-based methods.}
\label{tab:training_setup}
\scriptsize
\setlength{\tabcolsep}{5pt}
\renewcommand{\arraystretch}{1.05}
\begin{tabular}{
>{\raggedright\arraybackslash}p{0.17\textwidth}
>{\raggedright\arraybackslash}p{0.11\textwidth}
>{\raggedright\arraybackslash}p{0.21\textwidth}
>{\raggedright\arraybackslash}p{0.31\textwidth}}
\toprule
Setting & Value & Setting & Value \\
\midrule
Training steps & $1.0\times 10^6$ & Batch size / replay buffer & $256$ / $10^6$ \\
Discount / $n$-step return & $0.99$ / $3$ & Step-per-collect / epoch & $1024$ / $4096$ \\
Update-per-step & $0.5$ & Actor / critic learning rate & $10^{-4}$ / $2\times 10^{-5}$ \\
Weight decay / exploration noise & $10^{-4}$ / $0.05$ & Diffusion steps $K$ & $6$ \\
FNO width / retained modes & $48$ / $4$ & Spectral layers / activation & $1$ / Mish \\
Guidance scale $\eta$ & $0.5$ & BC weight schedule $\omega$ & $0.8 \rightarrow 0.1$ over $1.5\times 10^5$ updates \\
\bottomrule
\end{tabular}
\end{table*}

\textbf{Inference latency.} To make the deployment overhead explicit, we benchmarked the final Guided-DiffFNO policy on logged OfficeSmall evaluation states using a single NVIDIA RTX 3070 GPU. Including all six denoising steps and one critic-gradient computation per reverse step, the mean inference latency is $30.9$~ms per control decision for batch size $1$ (about $5.2$~ms per denoising step), which is small relative to the one-hour control interval used in this study.

\subsection{Evaluation Protocol and Metrics}
We report four primary metrics: average episode HVAC energy, average comfort violations, final test reward, and mean absolute action change. Table~\ref{tab:overall_results} reports the three-seed mean and standard deviation for all learning-based methods, where each seed reward is averaged over its last five evaluation checkpoints before seed aggregation. We additionally analyze cross-zone coordination, action and HVAC-power spectra, representative room trajectories, and the association between critic-estimated $Q(s,a)$ and Monte Carlo return. Unless noted otherwise, the plots in this section use three-seed statistics; only the local mechanism figures retain a canonical representative seed for interpretability.

These metrics capture complementary aspects of controller quality: energy and comfort violations reflect the primary task objective, final test reward summarizes the optimized RL objective, and mean absolute action change measures control smoothness. In the analysis, we therefore emphasize Pareto structure and consistency across scalar, spectral, and coordination results rather than drawing conclusions from any single metric alone.

\providecommand{\resultfigdir}{figures}

\section{Simulation Results}
\label{sec:results}

\begin{table*}[!t]
\centering
\caption{Overall comparison on the BEAR OfficeSmall-Hot\_Dry task. Learning-based methods report the three-seed mean $\pm$ standard deviation. Lower mean $|\Delta a|$ indicates smoother control.}
\label{tab:overall_results}
\footnotesize
\setlength{\tabcolsep}{4.5pt}
\renewcommand{\arraystretch}{1.12}
\begin{tabular}{lcccc}
\hline
Method & Energy (kWh) $\downarrow$ & Violations $\downarrow$ & Test Reward $\uparrow$ & Mean $|\Delta a|$ $\downarrow$ \\
\hline
Guided-DiffFNO & 871.5 $\pm$ 3.2 & \textbf{0.500 $\pm$ 0.030} & \textbf{-0.490 $\pm$ 0.018} & 0.0236 $\pm$ 0.0005 \\
DiffFNO w/o Guidance & 901.3 $\pm$ 19.9 & 0.886 $\pm$ 0.135 & -0.738 $\pm$ 0.086 & 0.0402 $\pm$ 0.0090 \\
DiffFNO w/o Residual \& Guidance & 871.1 $\pm$ 6.6 & 1.082 $\pm$ 0.338 & -0.844 $\pm$ 0.207 & 0.0357 $\pm$ 0.0109 \\
DiffFNO w/o Residual & \textbf{867.8 $\pm$ 5.3} & 0.614 $\pm$ 0.046 & -0.558 $\pm$ 0.029 & 0.0237 $\pm$ 0.0005 \\
Diffusion-MLP & 990.9 $\pm$ 38.4 & 1.441 $\pm$ 0.510 & -1.106 $\pm$ 0.321 & 0.0596 $\pm$ 0.0162 \\
MPC & 1004.9 & 1.780 & -1.325 & \textbf{0.0130} \\
SAC & 5980.4 $\pm$ 27.3 & 5.483 $\pm$ 0.046 & -5.055 $\pm$ 0.011 & 0.3430 $\pm$ 0.0412 \\
SAC+MPC & 1864.6 $\pm$ 74.6 & 3.431 $\pm$ 0.294 & -2.600 $\pm$ 0.222 & 0.1729 $\pm$ 0.0117 \\
\hline
\end{tabular}
\end{table*}

\subsection{Overall Energy--Comfort Trade-off}

Table~\ref{tab:overall_results} and Fig.~\ref{fig:energy_violations} show that Guided-DiffFNO and DiffFNO w/o Residual define the Pareto frontier. The latter attains the lowest raw energy ($867.8 \pm 5.3$~kWh), whereas Guided-DiffFNO attains the lowest comfort violations ($0.500 \pm 0.030$) and the best final reward ($-0.490 \pm 0.018$) with only a $3.8$~kWh/episode energy gap. On this benchmark, the full model is therefore the better-balanced controller when energy and thermal safety are considered jointly.

The gain from critic guidance is also clear. Relative to unguided DiffFNO, Guided-DiffFNO reduces energy by $3.3\%$, violations by $43.6\%$, and improves the final reward from $-0.738$ to $-0.490$, indicating that guidance changes the denoising trajectory in a practically meaningful way. Compared with Diffusion-MLP and MPC, the proposed method still achieves double-digit energy reductions together with markedly fewer violations, while the SAC-based baselines remain far from the relevant operating region.

\begin{figure}[!t]
\centering
\includegraphics[width=\linewidth]{\resultfigdir/compare_energy_violations.pdf}
\caption{Energy--comfort Pareto comparison. Points show three-seed means and error bars show one standard deviation. SAC and SAC+MPC are omitted for scale clarity.}
\label{fig:energy_violations}
\end{figure}

\subsection{Training Dynamics and Optimization Stability}

Figure~\ref{fig:reward_curves} shows that Guided-DiffFNO and DiffFNO w/o Residual form the strongest training trajectories, with Guided-DiffFNO reaching the best final reward and maintaining the narrowest terminal band. The unguided DiffFNO variants oscillate more and finish lower, while Diffusion-MLP and SAC+MPC do not sustain their stronger imitation-driven early phase. This suggests that structured diffusion plus critic guidance stabilizes long-horizon improvement more effectively than unguided denoising or conventional actor--critic training on this task.

\begin{figure}[!t]
\centering
\includegraphics[width=\linewidth]{\resultfigdir/compare_reward_curves.pdf}
\caption{Test reward curves during training. Curves show the three-seed mean and the shaded bands denote $\pm 1$ standard deviation.}
\label{fig:reward_curves}
\end{figure}

\subsection{Mechanism Validation of Critic Guidance}

Figure~\ref{fig:q_mc_return} tests whether the critic provides a meaningful guidance signal by comparing critic-estimated $Q(s,a)$ with Monte Carlo return on logged evaluation samples. The positive association ($r=0.619$, $\rho=0.638$, $N=504$) does not imply perfect calibration, but it does show that the critic preserves a useful ranking of better and worse actions. Together with the performance gap over unguided DiffFNO, this supports the interpretation of critic guidance as an effective inference-time refinement mechanism.

\begin{figure}[!t]
\centering
\includegraphics[width=0.88\linewidth]{\resultfigdir/critic_q_mc_return.pdf}
\caption{Relation between critic-estimated $Q(s,a)$ and Monte Carlo return for Guided-DiffFNO, computed from logged evaluation transitions of the representative seed.}
\label{fig:q_mc_return}
\end{figure}

\subsection{Action Quality and Physical Consistency}

\begin{figure}[!t]
\centering
\subfloat[Distribution of mean absolute action changes.\label{fig:action_smoothness_subfig}]{
\includegraphics[width=\linewidth]{\resultfigdir/compare_action_smoothness.pdf}}\\[2pt]
\subfloat[Welch power spectral density (PSD) of generated normalized action coordinates.\label{fig:action_psd_subfig}]{
\includegraphics[width=\linewidth]{\resultfigdir/action_psd_compare.pdf}}
\caption{Action-quality analysis for the learned controllers. Left: mean absolute action change across evaluation trajectories. Right: seed-averaged PSD of normalized action coordinates.}
\label{fig:action_quality}
\end{figure}

Beyond scalar reward and energy metrics, Guided-DiffFNO also produces smoother and more physically consistent control trajectories. In Fig.~\ref{fig:action_quality}, its mean absolute action change is $0.0236 \pm 0.0005$, essentially tied with DiffFNO w/o Residual and substantially lower than unguided DiffFNO ($0.0402 \pm 0.0090$), Diffusion-MLP ($0.0596 \pm 0.0162$), SAC+MPC ($0.1729 \pm 0.0117$), and SAC ($0.3430 \pm 0.0412$). The representative 72-hour coordination window in Fig.~\ref{fig:multizone_coordination} is consistent with this trend: Guided-DiffFNO shows strong but non-collapsed coordination across zones, with mean pairwise correlation $0.969$ and mean cross-zone standard deviation $0.015$, compared with $0.905/0.032$ for unguided DiffFNO and $0.560/0.061$ for Diffusion-MLP. The unguided FNO variant therefore already preserves substantial cross-zone coherence, but the guided full model remains more synchronized and less spatially dispersed, while the MLP baseline is clearly less coordinated.

\begin{figure}[!t]
\centering
\includegraphics[width=\linewidth]{\resultfigdir/multizone_action_coordination.pdf}
\caption{Representative 72-hour multi-zone action coordination for Guided-DiffFNO, DiffFNO w/o Guidance, and Diffusion-MLP.}
\label{fig:multizone_coordination}
\end{figure}

The spectral analysis leads to the same conclusion. Using a $2$ cycles/day cutoff on the discrete Welch spectrum, Guided-DiffFNO allocates $13.3\%$ of its action-spectrum mass to the high-frequency region, versus $27.9\%$ for unguided DiffFNO and $40.2\%$ for Diffusion-MLP. On the resulting HVAC power signal, the corresponding low-frequency mass ratios are $89.0\%$ for Guided-DiffFNO, $95.2\%$ for physical MPC, and $73.7\%$ for Diffusion-MLP. The learned policy therefore moves substantially closer to the low-frequency operating pattern of MPC while still achieving a better energy--comfort trade-off.

\begin{figure}[!t]
\centering
\includegraphics[width=\linewidth]{\resultfigdir/smalloffice_physical_psd_compare.pdf}
\caption{HVAC power spectral density for physical MPC, Guided-DiffFNO, and Diffusion-MLP. The dashed line marks the $2$ cycles/day cutoff.}
\label{fig:physical_psd_compare}
\end{figure}

The room-level trajectories in Fig.~\ref{fig:traj_control} provide a local illustration of the same pattern. Over the selected 72-hour window, Guided-DiffFNO leaves the comfort band only $3$ times, versus $12$ for unguided DiffFNO, and its room-wise action variation remains much smaller than that of Diffusion-MLP and SAC. The severe temperature excursions produced by SAC again indicate that standard point-estimate policies are poorly matched to this coupled control problem.

\begin{figure}[!t]
\centering
\subfloat[Temperature evolution for the selected room under Guided-DiffFNO and unguided DiffFNO.\label{fig:traj_temp_subfig}]{
\includegraphics[width=\linewidth]{\resultfigdir/temperature_trajectories_paper.pdf}}\\[2pt]
\subfloat[Control sequences for the same room generated by Guided-DiffFNO, Diffusion-MLP, and SAC.\label{fig:traj_control_subfig}]{
\includegraphics[width=\linewidth]{\resultfigdir/control_sequence_paper.pdf}}
\caption{Representative room-level trajectories over the selected 72-hour evaluation window.}
\label{fig:traj_control}
\end{figure}

\subsection{Ablation Analysis}

Figure~\ref{fig:ablation_heatmap} and Table~\ref{tab:overall_results} clarify the role of each component. The heatmap uses column-wise min--max normalization across the five ablation variants only; its smoothness axis is based on mean squared action increment, whereas absolute smoothness comparisons in the text use mean absolute action change.

Replacing the MLP denoiser with the FNO backbone already yields a clear gain in energy, comfort, and smoothness. Adding critic guidance further improves energy by $3.3\%$ and violations by $43.6\%$ relative to unguided DiffFNO, supporting the claim that guidance steers sampling toward safer and higher-value regions. Removing the residual branch keeps the controller smooth but slightly worsens comfort and reward, indicating that the residual pathway mainly restores local flexibility near comfort boundaries. Overall, the three components play complementary roles: FNO contributes structural smoothness, critic guidance contributes value alignment, and the residual branch improves local adaptability.

\begin{figure}[!t]
\centering
\includegraphics[width=\linewidth]{\resultfigdir/ablation_summary_heatmap.pdf}
\caption{Normalized ablation summary across energy, comfort, smoothness, and convergence. Each column is min--max normalized over the five ablation variants.}
\label{fig:ablation_heatmap}
\end{figure}

\subsection{Discussion}

From a control perspective, the main advantage of Guided-DiffFNO appears to be a better match between policy class and problem structure. Multi-zone HVAC control is coupled, safety-sensitive, and potentially multi-modal; SAC lacks expressive joint-action modeling, Diffusion-MLP lacks an explicit structural prior, and unguided diffusion lacks value-aware refinement. Guided-DiffFNO addresses these limitations jointly, which helps explain its more favorable compromise among energy, comfort, convergence, and smoothness.

An important nuance, however, is that the full model does not dominate every scalar criterion in isolation. Table~\ref{tab:overall_results} shows that DiffFNO w/o Residual attains slightly lower raw energy, and MPC remains the smoothest controller under the mean-$|\Delta a|$ metric. The benefit of the full model is therefore better understood as a stronger overall \emph{balance} across competing objectives. Consistent with the ablations and the spectral analyses, the FNO backbone mainly contributes a structured prior for coordinated low-frequency joint actions, whereas critic guidance improves value alignment and comfort compliance by steering sampling toward better action regions. The residual branch then preserves local flexibility near comfort boundaries, which helps explain why the full model achieves the most favorable overall energy--comfort compromise without uniformly dominating every baseline in every metric.

At the same time, the present study has clear boundaries. All experiments are conducted on one BEAR building configuration under one hot-dry climate profile and an hourly control interval. The conclusions are therefore about the value of structured generative policies in this benchmarked setting, not yet about universal superiority across building types, climates, or actuation frequencies. In addition, critic guidance should not be interpreted as a hard safety mechanism. Because the guidance signal is learned from reward, it can improve comfort compliance empirically, but it does not enforce explicit temperature constraints in the way a certified safety filter or a constrained optimizer would. The current framework also relies on MPC demonstrations for early stabilization, which is practical in many building studies but may become more restrictive when expert data are scarce or when the expert itself is misspecified.

These limitations suggest several concrete next steps. A first direction is broader external validation across buildings, climates, and occupancy patterns to determine when the structural prior transfers and when it should become graph- or topology-aware. A second direction is to integrate explicit safety mechanisms, for example constraint projection, shielding, or uncertainty-aware guidance, so that value-guided diffusion can be combined with formal comfort guarantees. A third direction is to study shorter control intervals and partial observability, where the interaction among diffusion depth, inference latency, and real-time deployment becomes more consequential. Even in its current form, however, the method illustrates a useful principle for building control: when the action space itself has meaningful physical structure, policy architectures that encode that structure can produce gains that are difficult to obtain through reward engineering or baseline scaling alone.

\section{Conclusion}
\label{sec:conclusion}

This paper proposed Guided-DiffFNO, a generative RL framework for multi-zone HVAC control that combines conditional diffusion, an FNO denoiser over the joint action vector, and inference-time critic guidance. On the BEAR \emph{OfficeSmall-Hot\_Dry} benchmark, the method achieves the best overall energy--comfort balance among the compared methods, with the lowest comfort violations, the highest final test reward, and smoother, more coordinated control trajectories than the alternative learned policies studied here.

The ablation and mechanism studies suggest that the method works because its components are complementary: the FNO backbone improves structural smoothness, critic guidance improves value alignment and comfort compliance, and the residual branch preserves local flexibility. At the same time, the present conclusions remain bounded by the current simulator benchmark and OfficeSmall configuration. Broader validation across additional buildings and climates, stronger safety mechanisms, longer-horizon generative planning, and eventual deployment under model mismatch remain important future directions.
\bibliographystyle{IEEEtran}
\bibliography{references.bib}

\end{document}
